\chapter{Consultas Avançadas, Junções e Subconsultas}
\label{cap:consultas_avancadas_joins_subconsultas}

Depois de dominar os comandos fundamentais da DML, o próximo passo é aprofundar a capacidade de análise por meio do comando \texttt{SELECT}. Em aplicações reais, raramente é suficiente listar registros de uma única tabela de forma isolada. Normalmente, é necessário combinar dados de várias tabelas, resumir informações, comparar grupos, localizar ausências e construir respostas mais ricas para problemas de negócio.

Esta aula amplia o estudo do \texttt{SELECT} com foco em:
\begin{itemize}
    \item funções de agregação;
    \item agrupamentos com \texttt{GROUP BY};
    \item filtros sobre grupos com \texttt{HAVING};
    \item subconsultas;
    \item junções entre tabelas com \texttt{JOIN}.
\end{itemize}

O domínio desses recursos marca a passagem da SQL operacional básica para a SQL analítica, aproximando o estudante da construção de relatórios, consultas gerenciais e exploração relacional de dados.

\section{Do SELECT simples ao SELECT analítico}

Em uma consulta básica, o usuário normalmente seleciona colunas e aplica filtros diretos. Exemplo:

\begin{sqlcode}
SELECT nome, curso
FROM alunos
WHERE status = 'ATIVO';
\end{sqlcode}

Esse tipo de comando é importante, mas ainda responde apenas a perguntas elementares. Em sistemas reais, surgem questões mais complexas, como:

\begin{itemize}
    \item quantos alunos existem por curso;
    \item qual a média de notas por disciplina;
    \item quais cursos ainda não possuem alunos;
    \item quais alunos têm nota acima da média;
    \item quais professores estão relacionados a turmas com maior número de matrículas.
\end{itemize}

Responder a perguntas assim exige raciocínio relacional mais elaborado.

\section{Banco de exemplo}

Para os exemplos desta aula, considere as seguintes tabelas:

\begin{sqlcode}
CREATE TABLE cursos (
    curso_id SERIAL PRIMARY KEY,
    nome VARCHAR(100) NOT NULL,
    coordenador VARCHAR(100)
);
\end{sqlcode}

\begin{sqlcode}
CREATE TABLE alunos (
    aluno_id SERIAL PRIMARY KEY,
    nome VARCHAR(100) NOT NULL,
    idade INTEGER,
    curso_id INTEGER,
    status VARCHAR(20) DEFAULT 'ATIVO',
    FOREIGN KEY (curso_id) REFERENCES cursos(curso_id)
);
\end{sqlcode}

\begin{sqlcode}
CREATE TABLE disciplinas (
    disciplina_id SERIAL PRIMARY KEY,
    nome VARCHAR(100) NOT NULL
);
\end{sqlcode}

\begin{sqlcode}
CREATE TABLE notas (
    nota_id SERIAL PRIMARY KEY,
    aluno_id INTEGER NOT NULL,
    disciplina_id INTEGER NOT NULL,
    nota NUMERIC(5,2) NOT NULL,
    FOREIGN KEY (aluno_id) REFERENCES alunos(aluno_id),
    FOREIGN KEY (disciplina_id) REFERENCES disciplinas(disciplina_id)
);
\end{sqlcode}

Essas tabelas serão suficientes para explorar agregações, subconsultas e junções em vários níveis.

\section{Funções de agregação}

As funções de agregação produzem um único valor a partir de várias linhas.

As principais são:
\begin{itemize}
    \item \texttt{COUNT()} -- conta ocorrências;
    \item \texttt{SUM()} -- soma valores;
    \item \texttt{AVG()} -- calcula média;
    \item \texttt{MAX()} -- retorna o maior valor;
    \item \texttt{MIN()} -- retorna o menor valor.
\end{itemize}

\subsection{Contagem total de alunos}

\begin{sqlcode}
SELECT COUNT(*) AS total_alunos
FROM alunos;
\end{sqlcode}

\subsection{Média de idade}

\begin{sqlcode}
SELECT AVG(idade) AS media_idade
FROM alunos;
\end{sqlcode}

\subsection{Maior e menor idade}

\begin{sqlcode}
SELECT MAX(idade) AS maior_idade,
       MIN(idade) AS menor_idade
FROM alunos;
\end{sqlcode}

\subsection{Média de notas}

\begin{sqlcode}
SELECT AVG(nota) AS media_geral
FROM notas;
\end{sqlcode}

\section{GROUP BY}

A cláusula \texttt{GROUP BY} organiza linhas em grupos e permite aplicar funções de agregação a cada grupo.

\subsection{Quantidade de alunos por curso}

\begin{sqlcode}
SELECT
    curso_id,
    COUNT(*) AS total_alunos
FROM alunos
GROUP BY curso_id;
\end{sqlcode}

\subsection{Média de idade por curso}

\begin{sqlcode}
SELECT
    curso_id,
    AVG(idade) AS media_idade
FROM alunos
GROUP BY curso_id;
\end{sqlcode}

\subsection{Média de nota por disciplina}

\begin{sqlcode}
SELECT
    disciplina_id,
    AVG(nota) AS media_disciplina
FROM notas
GROUP BY disciplina_id;
\end{sqlcode}

\begin{observacao}
Quando uma consulta usa agregação e, ao mesmo tempo, exibe outras colunas, essas colunas devem estar no \texttt{GROUP BY}, salvo quando forem resultado de função agregadora.
\end{observacao}

\section{HAVING}

A cláusula \texttt{HAVING} filtra grupos depois que eles já foram formados.

\subsection{Cursos com pelo menos três alunos}

\begin{sqlcode}
SELECT
    curso_id,
    COUNT(*) AS total_alunos
FROM alunos
GROUP BY curso_id
HAVING COUNT(*) >= 3;
\end{sqlcode}

\subsection{Disciplinas com média superior a 7}

\begin{sqlcode}
SELECT
    disciplina_id,
    AVG(nota) AS media_disciplina
FROM notas
GROUP BY disciplina_id
HAVING AVG(nota) > 7;
\end{sqlcode}

\begin{quadro_dica}
Uma forma prática de diferenciar \texttt{WHERE} e \texttt{HAVING} é lembrar que o primeiro filtra linhas antes do agrupamento, enquanto o segundo filtra grupos depois do agrupamento.
\end{quadro_dica}

\section{Subconsultas}

Uma subconsulta é uma consulta inserida dentro de outra. Ela é útil quando a lógica de decisão depende do resultado de uma consulta intermediária.

\subsection{Subconsulta com IN}

\begin{sqlcode}
SELECT nome
FROM alunos
WHERE curso_id IN (
    SELECT curso_id
    FROM cursos
    WHERE coordenador = 'Marcos Silva'
);
\end{sqlcode}

Nesse exemplo, a consulta externa retorna os nomes dos alunos vinculados a cursos coordenados por \texttt{Marcos Silva}.

\subsection{Subconsulta com comparação escalar}

\begin{sqlcode}
SELECT nome, idade
FROM alunos
WHERE idade > (
    SELECT AVG(idade)
    FROM alunos
);
\end{sqlcode}

Essa consulta retorna apenas os alunos cuja idade está acima da média geral.

\subsection{Subconsulta com EXISTS}

\begin{sqlcode}
SELECT nome
FROM alunos a
WHERE EXISTS (
    SELECT 1
    FROM notas n
    WHERE n.aluno_id = a.aluno_id
);
\end{sqlcode}

Aqui, são retornados apenas os alunos que possuem ao menos uma nota cadastrada.

\section{JOIN}

Junções são o mecanismo central de recomposição da informação em bancos normalizados. Como os dados estão distribuídos entre tabelas, o \texttt{JOIN} permite reuni-los logicamente em uma única resposta.

\section{INNER JOIN}

O \texttt{INNER JOIN} retorna apenas as linhas que possuem correspondência nas duas tabelas envolvidas.

\begin{sqlcode}
SELECT a.nome, c.nome AS curso
FROM alunos a
INNER JOIN cursos c
    ON a.curso_id = c.curso_id;
\end{sqlcode}

\subsection{INNER JOIN com três tabelas}

\begin{sqlcode}
SELECT a.nome AS aluno,
       d.nome AS disciplina,
       n.nota
FROM notas n
INNER JOIN alunos a
    ON n.aluno_id = a.aluno_id
INNER JOIN disciplinas d
    ON n.disciplina_id = d.disciplina_id;
\end{sqlcode}

Esse exemplo é muito próximo da realidade de relatórios acadêmicos.

\section{LEFT JOIN}

O \texttt{LEFT JOIN} retorna todas as linhas da tabela à esquerda, mesmo quando não há correspondência na tabela à direita.

\begin{sqlcode}
SELECT c.nome AS curso, a.nome AS aluno
FROM cursos c
LEFT JOIN alunos a
    ON c.curso_id = a.curso_id;
\end{sqlcode}

\subsection{Cursos sem alunos}

\begin{sqlcode}
SELECT c.nome
FROM cursos c
LEFT JOIN alunos a
    ON c.curso_id = a.curso_id
WHERE a.aluno_id IS NULL;
\end{sqlcode}

Esse tipo de consulta é útil para localizar ausências ou inconsistências.

\section{RIGHT JOIN}

O \texttt{RIGHT JOIN} devolve todas as linhas da tabela à direita e as correspondências da tabela à esquerda.

\begin{sqlcode}
SELECT a.nome, c.nome AS curso
FROM alunos a
RIGHT JOIN cursos c
    ON a.curso_id = c.curso_id;
\end{sqlcode}

Na prática, muitos desenvolvedores preferem reescrever esse raciocínio com \texttt{LEFT JOIN}, por clareza.

\section{FULL JOIN}

O \texttt{FULL JOIN} retorna todas as linhas quando existe correspondência em pelo menos um dos lados, preenchendo com \texttt{NULL} o lado ausente.

\begin{sqlcode}
SELECT a.nome, c.nome AS curso
FROM alunos a
FULL JOIN cursos c
    ON a.curso_id = c.curso_id;
\end{sqlcode}

Esse tipo de junção é útil para auditoria e análise de assimetrias entre conjuntos relacionados.

\section{Exemplo resolvido 1 -- Quantidade de alunos por curso}

Suponha que a coordenação deseje saber quantos alunos existem em cada curso.

\begin{sqlcode}
SELECT
    curso_id,
    COUNT(*) AS total_alunos
FROM alunos
GROUP BY curso_id;
\end{sqlcode}

\textbf{Interpretação:}
\begin{itemize}
    \item a tabela \texttt{alunos} é dividida em grupos por \texttt{curso\_id};
    \item a função \texttt{COUNT(*)} conta quantos alunos existem em cada grupo;
    \item o resultado pode ser usado em relatórios ou painéis.
\end{itemize}

\section{Exemplo resolvido 2 -- Alunos acima da média de idade}

Deseja-se localizar os alunos cuja idade está acima da média geral.

\begin{sqlcode}
SELECT nome, idade
FROM alunos
WHERE idade > (
    SELECT AVG(idade)
    FROM alunos
);
\end{sqlcode}

\textbf{Interpretação:}
\begin{itemize}
    \item a subconsulta calcula a média de idade;
    \item a consulta externa compara cada aluno com esse valor;
    \item somente alunos acima da média são retornados.
\end{itemize}

\section{Exemplo resolvido 3 -- Relatório de notas com nomes de alunos e disciplinas}

Como \texttt{notas} armazena apenas identificadores, é necessário usar junções para reconstruir a informação completa.

\begin{sqlcode}
SELECT a.nome AS aluno,
       d.nome AS disciplina,
       n.nota
FROM notas n
INNER JOIN alunos a
    ON n.aluno_id = a.aluno_id
INNER JOIN disciplinas d
    ON n.disciplina_id = d.disciplina_id;
\end{sqlcode}

\textbf{Interpretação:}
\begin{itemize}
    \item a consulta parte da tabela \texttt{notas};
    \item recupera o nome do aluno a partir de \texttt{aluno\_id};
    \item recupera o nome da disciplina a partir de \texttt{disciplina\_id};
    \item monta um relatório mais inteligível para usuários humanos.
\end{itemize}

\section{Exemplo resolvido 4 -- Cursos sem alunos}

A coordenação deseja identificar cursos ainda sem matrícula.

\begin{sqlcode}
SELECT c.nome
FROM cursos c
LEFT JOIN alunos a
    ON c.curso_id = a.curso_id
WHERE a.aluno_id IS NULL;
\end{sqlcode}

\textbf{Interpretação:}
\begin{itemize}
    \item todos os cursos são mantidos no resultado pelo \texttt{LEFT JOIN};
    \item quando não há aluno correspondente, as colunas da tabela \texttt{alunos} ficam nulas;
    \item o \texttt{WHERE a.aluno\_id IS NULL} filtra justamente esses casos.
\end{itemize}

\section{Exemplo resolvido 5 -- Disciplinas com média superior a 7}

\begin{sqlcode}
SELECT
    disciplina_id,
    AVG(nota) AS media_disciplina
FROM notas
GROUP BY disciplina_id
HAVING AVG(nota) > 7;
\end{sqlcode}

\textbf{Interpretação:}
\begin{itemize}
    \item as notas são agrupadas por disciplina;
    \item calcula-se a média de cada grupo;
    \item o \texttt{HAVING} retém apenas os grupos com média superior a 7.
\end{itemize}

\section{Erro comum 1 -- Esquecer a condição de junção}

Exemplo incorreto:

\begin{sqlcode}
SELECT a.nome, c.nome
FROM alunos a, cursos c;
\end{sqlcode}

Sem a condição que relaciona as tabelas, o resultado pode virar um produto cartesiano, multiplicando artificialmente o número de linhas.

\section{Erro comum 2 -- Usar WHERE em vez de HAVING}

Exemplo inadequado:

\begin{sqlcode}
SELECT curso_id, COUNT(*)
FROM alunos
WHERE COUNT(*) >= 3
GROUP BY curso_id;
\end{sqlcode}

Isso não funciona corretamente porque \texttt{COUNT(*)} é uma agregação. O filtro sobre grupos deve ser feito em \texttt{HAVING}.

\section{Erro comum 3 -- Misturar colunas agregadas e não agregadas sem GROUP BY}

Exemplo incorreto:

\begin{sqlcode}
SELECT curso_id, nome, COUNT(*)
FROM alunos;
\end{sqlcode}

Sem agrupar corretamente, a consulta torna-se semanticamente inválida.

\section{Erro comum 4 -- Usar subconsulta quando uma junção seria mais clara}

Subconsultas são úteis, mas nem sempre são a solução mais legível. Em muitos casos, uma junção expressa melhor a lógica relacional.

\section{Boas práticas}

\begin{itemize}
    \item usar aliases curtos e consistentes;
    \item construir consultas em etapas;
    \item validar a lógica com conjuntos pequenos antes de aplicar em tabelas maiores;
    \item preferir nomes de colunas explícitos em vez de \texttt{SELECT *};
    \item usar \texttt{LEFT JOIN} com atenção ao significado de \texttt{NULL};
    \item revisar se o problema é de linha, grupo ou relação entre tabelas.
\end{itemize}

\begin{quadro_dica}
Uma estratégia muito útil é montar a consulta por camadas: primeiro as tabelas e junções, depois os filtros, em seguida os agrupamentos e, por fim, a ordenação.
\end{quadro_dica}

\section{Minicaso integrador -- Relatório acadêmico de acompanhamento}

Considere que a coordenação deseja um relatório com:

\begin{enumerate}
    \item a quantidade de alunos por curso;
    \item os cursos sem alunos;
    \item os alunos com idade acima da média;
    \item a listagem de notas com aluno e disciplina;
    \item as disciplinas cuja média de notas supera 7.
\end{enumerate}

Uma sequência de consultas possível é:

\subsection*{Passo 1 -- Quantidade de alunos por curso}

\begin{sqlcode}
SELECT
    curso_id,
    COUNT(*) AS total_alunos
FROM alunos
GROUP BY curso_id;
\end{sqlcode}

\subsection*{Passo 2 -- Cursos sem alunos}

\begin{sqlcode}
SELECT c.nome
FROM cursos c
LEFT JOIN alunos a
    ON c.curso_id = a.curso_id
WHERE a.aluno_id IS NULL;
\end{sqlcode}

\subsection*{Passo 3 -- Alunos acima da média de idade}

\begin{sqlcode}
SELECT nome, idade
FROM alunos
WHERE idade > (
    SELECT AVG(idade)
    FROM alunos
);
\end{sqlcode}

\subsection*{Passo 4 -- Relatório de notas}

\begin{sqlcode}
SELECT a.nome AS aluno,
       d.nome AS disciplina,
       n.nota
FROM notas n
INNER JOIN alunos a
    ON n.aluno_id = a.aluno_id
INNER JOIN disciplinas d
    ON n.disciplina_id = d.disciplina_id;
\end{sqlcode}

\subsection*{Passo 5 -- Disciplinas com média maior que 7}

\begin{sqlcode}
SELECT
    disciplina_id,
    AVG(nota) AS media_disciplina
FROM notas
GROUP BY disciplina_id
HAVING AVG(nota) > 7;
\end{sqlcode}

\textbf{Comentário didático:}
Esse minicaso mostra que consultas mais ricas raramente dependem de um único recurso. Em geral, elas combinam agregação, junção e subconsulta para responder a perguntas analíticas completas.

\section{Exercícios básicos}

\begin{enumerate}
    \item Conte o número total de alunos.
    \item Calcule a média de idade dos alunos.
    \item Liste a quantidade de alunos por curso.
    \item Liste apenas os cursos com pelo menos dois alunos.
    \item Retorne os alunos cuja idade esteja acima da média geral.
\end{enumerate}

\section{Exercícios intermediários}

\begin{enumerate}
    \item Faça uma consulta com \texttt{INNER JOIN} entre \texttt{alunos} e \texttt{cursos}.
    \item Faça uma consulta com \texttt{LEFT JOIN} para localizar cursos sem alunos.
    \item Monte uma consulta que traga aluno, disciplina e nota.
    \item Escreva uma subconsulta que filtre alunos vinculados a cursos coordenados por um professor específico.
    \item Explique a diferença entre \texttt{WHERE} e \texttt{HAVING}.
\end{enumerate}

\section{Exercícios avançados}

\begin{enumerate}
    \item Crie uma consulta que combine \texttt{JOIN} e agregação.
    \item Produza uma consulta com \texttt{EXISTS} para localizar alunos com nota cadastrada.
    \item Elabore um exemplo de \texttt{FULL JOIN} e explique por que ele pode ser útil.
    \item Reescreva uma subconsulta usando junção, quando possível.
\end{enumerate}

\section{Exercício integrador}

Considere um sistema universitário com as tabelas \texttt{alunos}, \texttt{cursos}, \texttt{professores}, \texttt{disciplinas} e \texttt{notas}. Elabore um conjunto de consultas que responda:

\begin{enumerate}
    \item quantos alunos existem por curso;
    \item quais cursos não possuem alunos;
    \item quais alunos têm nota acima de 8;
    \item quais alunos estão acima da média de idade;
    \item quais disciplinas possuem média superior a 7;
    \item quais alunos possuem ao menos uma nota registrada.
\end{enumerate}

\section{Síntese da aula}

Esta aula aprofundou o uso analítico da SQL ao estudar funções de agregação, agrupamentos, filtros de grupo, subconsultas e junções. Com exemplos resolvidos, erros comuns, exercícios progressivos e um minicaso integrador, o estudante passa a compreender como o \texttt{SELECT} pode ser usado para responder perguntas mais sofisticadas sobre dados relacionais. A próxima aula aprofundará o refinamento de consultas com cláusulas adicionais e introduzirá o uso de \texttt{VIEW}.